\documentclass[11pt,a4paper]{ctexart}

\usepackage[margin=2.5cm]{geometry}
\usepackage{enumitem}
\usepackage{amsmath,amssymb}
\usepackage[colorlinks=true,linkcolor=blue,urlcolor=blue]{hyperref}

\setlist[enumerate]{itemsep=2pt,topsep=3pt}
\linespread{1.15}

\title{\bfseries 大学物理竞赛\\[0.3em]（University Physics Competition）\\[0.3em]参赛规则}
\author{}
\date{}

\begin{document}

\maketitle

\vspace{-2em}


\vspace{1em}

\section{指导教师}

要参加大学物理竞赛，每支参赛队必须由所在院校的一名教师担任该队的指导教师（Faculty
Sponsor）。指导教师的职责是：确认参赛学生确实在本校注册就读、确认他们符合参赛资格，并保证本队将遵守竞赛规则。指导教师不必在整整 48 小时的竞赛过程中全程现场监督参赛队。

竞赛开始之前，指导教师可以帮助参赛队做准备，例如与学生讨论解题策略、推荐有用的资料等。竞赛开始、赛题公布之后，指导教师不得再做出任何贡献。竞赛开始之后，学生不得与指导教师或队外的任何人讨论。一名教师可以担任本竞赛中多支参赛队的指导教师，对其可指导的参赛队数量没有任何限制。

\section{竞赛报名}

所有参赛队都必须在本网站上正式注册，并在竞赛开始前至少六小时缴纳报名费。提前报名（可享受优惠报名费）必须在竞赛开始前至少 7 天完成。参赛队将在提交报名后的两个工作日内、且不迟于竞赛开始前 3 小时，收到一封分配队伍编号（Team
Number）的电子邮件。报名时必须列出队内学生的姓名。如果指导教师需要更改队内学生姓名，请发送邮件至
\href{mailto:director@uphysicsc.com}{\mbox{director@uphysicsc.com}}。竞赛开始之后，队员不得更换。一所院校可以派出的参赛队数量没有限制。

如果某支队伍无法参赛，只要在竞赛开始前提出退款申请，我们将退还报名费。竞赛一旦开始，恕不退款。

\section{参赛队}

每支参赛队最多由三名学生组成（因此两人甚至一人组队也是允许的）。每名学生只能参加一支队伍。队员在竞赛期间必须是在校注册的学生，但不要求是全时制学生。队员必须与指导教师及其他队员注册于同一所院校。该院校可以是任何高中后（post-secondary）院校、大学或学院。队员不得拥有任何 STEM 领域（科学、技术、工程、数学）的学士学位。

\section{赛题}

竞赛开始时，赛题将在本网站 \url{https://www.uphysicsc.com} 上公布。届时将提供两道赛题。每支参赛队只能就两道赛题中的一道提交解答。

\section{解答论文}

所有参赛队都必须用英文撰写一篇正式论文，在其中说明自己所做的工作。

论文必须在每一页的顶部标明队伍编号，但论文不得包含学生姓名、所在院校或任何其他可识别身份的信息。论文的第一页只应包含四项内容：(1) 队伍编号；(2) 所选择的赛题（“A”或“B”）；(3) 论文标题；(4) 一篇不超过
300 词的摘要（summary），该摘要应为一段提供所做工作关键细节、所用模型的假设、对所提问题的回答、最终数值、结果与结论的摘要性文字（abstract）。位于这一页之后的论文正文，长度不限，以足以报告本队工作为准。

竞赛期间，参赛队可以使用书籍、期刊、计算机、互联网、自己编写的程序或任何其他非生命的资源，但不得与队外的任何人商讨。指导教师或除这三名学生以外的其他任何人，都不得做出任何贡献。虽然赛题的本意是进行理论分析，参赛队也可以开展相关实验，并在论文中呈现由此得到的数据。

每篇论文都必须包含所用参考文献列表，并在正文中给出对应的引用标注。论文可以使用已有的算法和计算工具（不仅包括免费公开的工具，也包括 Matlab 或 Mathematica 等系统中的工具），前提是这些工具被正确地引用和标注，并且其方法在论文中被清楚地说明。论文可以包含公式、曲线图、插图和表格。数值必须给出适当的有效数字，并清楚标明单位。曲线图必须有清楚标注的坐标轴。所有变量都必须清楚定义，并在全文保持一致。任何自行编写的计算机程序可以作为附录收入论文，但所有算法、方法和结果都必须在论文正文中加以说明，才能在评审时获得考虑。

每篇论文至少必须包含：(1) 以本队自己的理解对问题的重述与澄清；(2) 对所有假设与近似及其理由的说明；(3) 对所进行的全部工作的论证；(4) 对所采用方法的优点与不足的讨论。

每支参赛队只能就两道赛题中的一道提交论文。如果一支队伍提交了两篇论文、分别对应两道赛题，则只评审问题
A 的那篇论文。

\subsection*{生成式人工智能（Generative AI）}

鼓励参赛队不使用生成式 AI 工具（例如 ChatGPT、Jasper、Anyword、Rytr、Bard 或类似系统）来撰写论文，因为
AI 工具生成的内容存在若干问题，在物理这样的技术学科中尤其如此。人们发现，AI
工具会产生幻觉，即生成虚假的内容；也发现有 AI 工具进行剽窃，即未经引用和归属说明就照搬其他来源的内容。包含剽窃或虚假内容的论文将受到严厉惩处。因此，务必彻底核查任何由
AI 生成的内容。

如果参赛队在撰写论文时使用了生成式 AI 工具，则必须就此保持透明，清楚地说明使用了哪些工具以及以何种方式使用。论文末尾必须包含一节“人工智能使用报告”（Report
on the Use of AI），说明：(1) 使用了哪个系统；(2) 查询内容（输入该 AI 工具的确切措辞）；(3)
完整输出（该 AI 工具完整输出的副本）。

\section{提交解答}

论文可以 pdf、doc 或 docx 文件形式提交。每支参赛队必须将论文通过电子邮件发送至
\href{mailto:solutions@uphysicsc.com}{\mbox{solutions@uphysicsc.com}}，且论文必须在竞赛结束之前送达。竞赛在开始后恰好
48 小时结束。如果某支队伍的论文未在该截止时间之前通过电子邮件送达，该队将不被纳入奖牌排名。

如果一支已正确报名的参赛队所提交的论文在 48 小时截止时间之后略微迟到才被收到，评审可以自行决定将这篇论文评为“成功参赛者”（accomplished
competitor）。但是，要获得奖牌排名，论文必须在 48 小时截止时间之前送达。

论文的文件名应为队伍编号后接赛题字母，例如 778A.pdf 或 832B.doc。队伍编号应在邮件的主题栏和邮件正文中都写明。整篇论文必须作为邮件附件提交。不要提交指向
Google Docs 等云服务上论文的链接。

\section{竞赛结果}

竞赛结束后，所有论文都将被评审并排名。最佳的大约 2\% 的论文将被评为金牌获得者；接下来的
15\%--20\% 的论文将被评为银牌获得者；再接下来的 25\%--30\% 的论文将被评为铜牌获得者；其余
50\%--60\% 的论文将被评为成功参赛者。竞赛结果公布后，每支成功参赛的队伍都将获得一份可下载的 pdf
证书。

\section{竞赛主任}

大学物理竞赛主任（Director）是所有规则与政策的最终裁决者，并可自行决定取消任何在其单独裁量下违反了竞赛规则的参赛队的资格。

\end{document}
